% Основной текст. Подключается из stage1_report.tex.
\section{Постановка задачи и статус согласования}
\subsection{Тема и цель работы}
Предлагаемая предметная область --- запись клиентов на услуги салона красоты. Тема курсовой работы: «Информационная система онлайн-записи клиентов в салон красоты».

Цель системы --- предоставить клиенту возможность самостоятельно выбрать услугу, мастера, дату и свободное время, а администратору --- вести услуги, сведения о мастерах, их рабочие интервалы и общий журнал записей. Система должна учитывать продолжительность услуги и предотвращать пересечение записей одного мастера.

В настоящем отчёте представлены описание предметной области, обоснование автоматизации, требования, модели и описания основных прецедентов, а также предлагаемая архитектура. Реализация прикладного кода, физическая схема базы данных и развёртывание относятся к последующим этапам.

\subsection{Вопросы, подлежащие согласованию с преподавателем}
Предметная область, набор основных прецедентов и технологии в этом отчёте являются \textbf{предложением для согласования}. Факт их утверждения преподавателем не предполагается.
\begin{enumerate}
\item Утвердить тему и границы системы: один салон, две роли, индивидуальные услуги по предварительной записи.
\item Утвердить основные прецеденты: поиск свободного времени и запись; просмотр и отмена своих записей; управление услугами, мастерами, рабочими интервалами и журналом записей.
\item Утвердить стек Vue, Spring Boot со Spring MVC и PostgreSQL.
\item Уточнить доступную версию Java и PostgreSQL, правила запуска процессов, выделение порта и способ доступа к демонстрации на сервере \texttt{helios}.
\end{enumerate}
Дата и результат согласования: \textit{заполняются после обсуждения с преподавателем}.

\section{Описание предметной области}
\subsection{Деятельность салона и участники процесса}
Салон красоты оказывает услуги по предварительной записи: например, стрижку, укладку и уход за лицом. Каждая услуга имеет название, описание, стоимость и фиксированную продолжительность в минутах. Приведённые услуги являются примерами, а не обязательным каталогом.

Услуги выполняют мастера. Один мастер может выполнять несколько услуг, а одну услугу могут выполнять несколько мастеров. Продолжительность и стоимость услуги в базовой версии одинаковы для всех выполняющих её мастеров. В каждый момент времени мастер обслуживает не более одного клиента.

Для каждого мастера администратор задаёт рабочие интервалы на конкретные даты. Например, работа с 09:00 до 13:00 и с 14:00 до 18:00 описывается двумя интервалами. Промежуток между ними недоступен для записи. Выходной день задаётся отсутствием рабочих интервалов. Такой подход позволяет учитывать перерывы и исключения без отдельного механизма повторяющихся графиков.

Клиент выбирает услугу, подходящего мастера и дату, после чего видит допустимые моменты начала посещения. При выборе времени система показывает также время окончания и стоимость. Запись закрепляется за клиентом только после успешного сохранения на сервере.

Администратор ведёт справочники и графики, просматривает записи по датам, клиентам, мастерам и статусам, а при необходимости отменяет запись. Мастер в базовой версии является объектом учёта и не получает отдельного личного кабинета.

\subsection{Процесс без информационной системы}
В рассматриваемой исходной модели клиент обращается по телефону или через сообщения. Администратор уточняет желаемую услугу, находит мастера, вручную сравнивает продолжительность услуги с графиком и журналом записей, затем согласует время с клиентом. При отмене администратор вручную освобождает соответствующий интервал.

Такая организация создаёт риск двойной записи, ошибок при определении времени окончания, предложения времени в перерыв и потери сведений об отмене. Клиент не видит актуальную доступность самостоятельно и зависит от ответа администратора. Это аналитическая модель исходного процесса; обследование конкретного действующего салона в рамках данного отчёта не проводилось.

\subsection{Границы базовой версии}
\begin{itemize}
\item Учитывается один салон с одним настроенным часовым поясом.
\item Одна запись содержит одного клиента, одного мастера и одну услугу.
\item Записи не переходят через полночь; услуга целиком помещается в один рабочий интервал.
\item Доступное время предлагается на ближайшие 30 календарных дней, включая текущую дату; прошедшее время исключается.
\item Начало записи выбирается с шагом 15 минут. Продолжительность услуги --- положительное целое число минут и не обязана быть кратной этому шагу.
\item Подтверждение автоматическое: отдельное одобрение администратором и временное удержание выбранного времени не предусмотрены.
\item Клиент может отменить запись до её начала. Администратор также отменяет только будущие записи, указывая причину.
\end{itemize}
В базовую версию не входят онлайн-оплата, рассылки, интеграции с внешними календарями, программы лояльности, складской учёт, несколько филиалов, групповые услуги, учёт кабинетов и оборудования. Перенос выполняется отменой старой записи и созданием новой; сохранение прежнего времени при таком переносе не гарантируется. Параллельные записи одного клиента к разным мастерам дополнительно не ограничиваются: обязательный ресурсный конфликт определяется по мастеру.

\subsection{Основные понятия и сущности}
\begin{longtable}{L{3.3cm}L{11.5cm}}
\toprule Сущность & Содержание и основные сведения \\ \midrule
\endhead
Пользователь & Идентификатор, имя, адрес электронной почты для входа, хеш пароля, роль: клиент или администратор. Открытая регистрация создаёт только клиента. \\
Мастер & Идентификатор, имя, краткое описание квалификации, признак активности. Не обязан иметь учётную запись пользователя. \\
Услуга & Идентификатор, название, описание, длительность в минутах, стоимость, признак активности. \\
Мастер--услуга & Связь мастера с услугой, которую он вправе выполнять; пара идентификаторов уникальна. \\
Рабочее время мастера & Идентификатор, мастер, календарная дата, начало и окончание рабочего интервала. Интервалы одного мастера не пересекаются. \\
Запись & Идентификатор, клиент, мастер, услуга, начало и окончание, зафиксированные длительность и стоимость, статус, дата создания; для отмены --- время, инициатор и причина. \\
\bottomrule
\end{longtable}
Один клиент имеет много записей; один мастер имеет много записей и рабочих интервалов; одна услуга может встречаться во многих записях. Связь мастеров и услуг имеет тип «многие ко многим» и представлена отдельной сущностью. Каждая запись ссылается ровно на одного клиента, мастера и услугу.

Свободный интервал не является самостоятельной хранимой сущностью: он вычисляется из услуги, графика и действующих записей. Длительность и стоимость фиксируются при записи, поэтому последующее изменение справочника не изменяет уже согласованные условия посещения. Сущности, на которые имеются ссылки, не удаляются физически: услуги и мастера деактивируются.

\section{Назначение и ожидаемые результаты автоматизации}
Система нужна для ведения единого актуального расписания и самостоятельной записи клиентов. Она решает следующие задачи:
\begin{enumerate}
\item Предоставляет клиенту каталог услуг и список мастеров, способных выполнить выбранную услугу.
\item Автоматически вычисляет допустимые моменты начала с учётом полной длительности услуги, рабочих интервалов и существующих записей.
\item Обеспечивает согласованное сохранение записей при одновременном обращении нескольких клиентов.
\item Даёт клиенту доступ к собственным будущим и прошедшим записям и возможность отмены будущей записи.
\item Позволяет администратору централизованно поддерживать справочники и рабочее расписание, контролировать журнал записей.
\end{enumerate}
Ожидаемый результат --- снижение числа ручных операций и устранение двойных записей, создаваемых через систему. Количественная экономия времени не заявляется без измерений; достижение цели оценивается по сформулированным требованиям.

\section{Требования к системе}
\subsection{Функциональные требования}
\begin{longtable}{L{1.6cm}L{13.2cm}}
\toprule Код & Требование \\ \midrule
\endhead
ФТ-01 & Система должна поддерживать регистрацию клиента, вход и выход. Почта уникальна; пароль не хранится в открытом виде. Создание администратора выполняется при начальной настройке. \\
ФТ-02 & Клиент должен просматривать активные услуги с описанием, стоимостью и длительностью, а также выбирать только активных мастеров, выполняющих услугу. \\
ФТ-03 & После выбора услуги, мастера и даты система должна вычислять доступные начала с учётом длительности услуги, рабочего времени мастера и существующих действующих записей и показывать отсутствие свободного времени явно. \\
ФТ-04 & Авторизованный клиент должен создавать запись после просмотра услуги, мастера, начала, окончания и стоимости. Сервер определяет клиента из текущей сессии. \\
ФТ-05 & Система должна повторно проверять доступность при сохранении и предотвращать пересекающиеся действующие записи одного мастера, включая конкурентные запросы. \\
ФТ-06 & Клиент должен видеть только собственные записи, включая отменённые и прошедшие, и отменять собственную действующую запись до её начала. \\
ФТ-07 & Администратор должен создавать и редактировать услуги и мастеров, управлять активностью и связями «мастер--услуга» с соблюдением ограничений будущих записей. \\
ФТ-08 & Администратор должен задавать, изменять и удалять рабочие интервалы на конкретные даты. Система должна отклонять пересечения рабочих интервалов и изменения, нарушающие будущие записи. \\
ФТ-09 & Администратор должен просматривать общий журнал с фильтрами по дате, мастеру, клиенту и статусу и отменять будущую запись с указанием причины. \\
ФТ-10 & Система должна сохранять историю отмен и согласованные длительность и стоимость записи; изменение справочника не должно переписывать эти условия. \\
ФТ-11 & При ошибке данных, недостатке прав, конфликте времени или недоступности сервера система должна сообщать причину на понятном русском языке и не показывать ложное подтверждение. \\
\bottomrule
\end{longtable}

\subsection{Нефункциональные требования}
Приведённые показатели являются целевыми критериями для будущей проверки, а не результатами уже выполненных испытаний.
\begin{longtable}{L{1.6cm}L{13.2cm}}
\toprule Код & Требование и способ проверки \\ \midrule
\endhead
НФТ-01 & Целостность: создание записи и отмена выполняются атомарно. При двух конкурентных запросах на один интервал сохраняется ровно одна запись; проверяется интеграционным сценарием с общей БД. \\
НФТ-02 & Разграничение доступа: сервер проверяет роль и владельца объекта на каждом защищённом запросе. Клиент не получает чужие записи и административные функции даже при подмене идентификаторов. \\
НФТ-03 & Защита данных: пароли хранятся как стойкие адаптивные хеши с солью; сессии используют защищённые настройки cookie, изменяющие запросы защищаются от CSRF. Пароли, cookie и полные реквизиты подключения не попадают в журналы. \\
НФТ-04 & Производительность: целевой 95-й процентиль времени серверного расчёта доступности --- не более 2 секунд при 20 одновременных клиентах, 20 мастерах и 10\,000 записях. Проверяется на согласованной конфигурации \texttt{helios}; сетевое время учитывается отдельно. \\
НФТ-05 & Удобство: интерфейс на русском языке, пригоден для экрана шириной от 360 пикселей; выбор выполняется в порядке «услуга --- мастер --- дата --- время --- подтверждение». Кнопки доступны с клавиатуры, состояние не передаётся только цветом. \\
НФТ-06 & Надёжность: подтверждение выдаётся только после фиксации транзакции. После перезапуска приложения сохранённые записи доступны. При неопределённом результате сетевого запроса клиент может проверить список своих записей. \\
НФТ-07 & Развёртывание: серверная часть запускается на \texttt{helios} в рамках разрешений пользователя; браузер обращается к экземпляру на этом сервере. Наличие Docker, прав администратора и отдельного Node.js-процесса не требуется. \\
НФТ-08 & Сопровождаемость: интерфейс, бизнес-правила и доступ к данным разделены; конфигурация окружения вынесена из исходного кода. Расчёт интервалов и конкурентная запись проверяются автоматически на этапе реализации. \\
НФТ-09 & Единое время: даты отображаются в настроенном часовом поясе салона, моменты записей хранятся однозначно с учётом смещения. Авторитетным источником текущего времени служит сервер. \\
\bottomrule
\end{longtable}

\section{Модели и описания основных прецедентов}
\subsection{Акторы и граница системы}
\textbf{Клиент} выбирает услуги и управляет собственными записями. \textbf{Администратор} поддерживает справочники и графики, управляет общим журналом. Неавторизованный посетитель рассматривается как клиент до входа, а не как отдельная бизнес-роль. База данных и внутренний расчёт времени не являются внешними акторами.

На рисунках ассоциации показывают участие актора, а пунктирная стрелка \texttt{include} направлена к обязательному включаемому прецеденту. Вход в систему является отдельным прецедентом и предусловием защищённых операций.

\begin{figure}[htbp]
\centering
\begin{tikzpicture}[font=\small,
 usecase/.style={ellipse,draw,align=center,text width=5.1cm,minimum height=0.85cm},
 inc/.style={dashed,-{Stealth}}]
\node (actor) at (0,-0.5) {Клиент};
\draw (0,0.8) circle (0.16);
\draw (0,0.64)--(0,0.2) (-0.26,0.47)--(0.26,0.47) (0,0.2)--(-0.23,-0.12) (0,0.2)--(0.23,-0.12);
\node[usecase] (u1) at (5,2.6) {П1. Регистрация и вход};
\node[usecase] (u2) at (5,1.3) {П2. Поиск свободного времени};
\node[usecase] (u3) at (5,0) {П3. Создание записи};
\node[usecase] (u4) at (5,-1.3) {П4. Просмотр своих записей};
\node[usecase] (u5) at (5,-2.6) {П5. Отмена своей записи};
\node[usecase] (u9) at (5,-4.3) {П9. Проверка доступности};
\foreach \u in {u1,u2,u3,u4,u5} \draw (actor.east)--(\u.west);
\draw[inc] (u3.east) -- ++(1.1,0) |- node[pos=0.3,right]{\scriptsize\texttt{<<include>>}} (u9.east);
\begin{scope}[on background layer]
\node[draw,fit=(u1)(u9),inner xsep=1.65cm,inner ysep=0.5cm,label=above:{Система онлайн-записи: клиент}] {};
\end{scope}
\end{tikzpicture}
\caption{Модель клиентских прецедентов}
\end{figure}

\begin{figure}[htbp]
\centering
\begin{tikzpicture}[font=\small,
 usecase/.style={ellipse,draw,align=center,text width=6.1cm,minimum height=1cm}]
\node (actor) at (0,-0.5) {Администратор};
\draw (0,0.9) circle (0.16);
\draw (0,0.74)--(0,0.3) (-0.26,0.55)--(0.26,0.55) (0,0.3)--(-0.23,0.02) (0,0.3)--(0.23,0.02);
\node[usecase] (u1) at (6,2.3) {П1. Вход};
\node[usecase] (u6) at (6,0.8) {П6. Управление услугами,\\мастерами и их связями};
\node[usecase] (u7) at (6,-0.8) {П7. Управление рабочим временем};
\node[usecase] (u8) at (6,-2.3) {П8. Работа с журналом записей};
\foreach \u in {u1,u6,u7,u8} \draw (actor.east)--(\u.west);
\begin{scope}[on background layer]
\node[draw,fit=(u1)(u8),inner sep=0.45cm,label=above:{Система онлайн-записи: администратор}] {};
\end{scope}
\end{tikzpicture}
\caption{Модель административных прецедентов}
\end{figure}
\clearpage

\subsection{П1. Регистрация и вход}
\ucfield{Акторы}{Клиент; администратор --- только вход и выход.}
\ucfield{Предусловия}{Для входа существует учётная запись; для регистрации клиент не авторизован.}
\ucfield{Триггер}{Пользователь открывает форму регистрации или входа.}
\ucfield{Основной сценарий}{Клиент вводит имя, почту и пароль. Система проверяет формат и уникальность почты, сохраняет учётную запись с ролью клиента. Пользователь вводит почту и пароль на форме входа; система проверяет их и открывает доступ согласно роли.}
\ucfield{Альтернативы}{При занятой почте или некорректных данных регистрация отклоняется. При неверных реквизитах входа сессия не создаётся. При выходе сессия прекращается.}
\ucfield{Постусловия}{Создана учётная запись либо открыта авторизованная сессия; при отказе права доступа не изменились.}

\subsection{П2. Поиск свободного времени}
\ucfield{Актор}{Клиент.}
\ucfield{Предусловия}{В системе опубликованы активные услуги; вход необязателен.}
\ucfield{Триггер}{Клиент начинает подбор времени посещения.}
\ucfield{Основной сценарий}{Клиент выбирает услугу. Система показывает выполняющих её активных мастеров. Клиент выбирает мастера и дату. Система получает график и действующие записи, рассчитывает варианты начала, показывает их вместе с длительностью и стоимостью услуги.}
\ucfield{Альтернативы}{Если мастер не работает, свободного времени нет или для услуги нет мастеров, система сообщает об этом и предлагает изменить выбор. Дата вне разрешённого горизонта отклоняется.}
\ucfield{Постусловия}{Клиент получил варианты; записи в БД не созданы, время не удерживается.}

\subsection{П3. Создание записи}
\ucfield{Актор}{Клиент.}
\ucfield{Предусловия}{Клиент вошёл в систему; выбраны услуга, мастер, дата и предложенное время.}
\ucfield{Триггер}{Клиент нажимает кнопку подтверждения.}
\ucfield{Основной сценарий}{
\begin{enumerate}
\item Система показывает итоговые условия: услугу, мастера, начало, окончание и стоимость.
\item Клиент подтверждает условия.
\item Сервер начинает транзакцию, получает необходимые блокировки и выполняет П9.
\item При успехе сервер сохраняет действующую запись с зафиксированными условиями и завершает транзакцию.
\item Клиент получает номер записи и видит её в своём списке.
\end{enumerate}}
\ucfield{Альтернативы}{При конфликте времени запись не создаётся и список времени обновляется. При изменении стоимости или длительности требуется повторное подтверждение актуальных условий. При истечении сессии требуется вход. При ошибке сохранения выполняется откат. Если ответ потерян после фиксации, клиент проверяет свои записи перед повторением; повторный запрос на тот же занятый интервал не создаёт дубликат.}
\ucfield{Постусловия}{Создана ровно одна согласованная запись либо новая запись отсутствует; ранее сохранённые записи не изменены.}

\subsection{П4. Просмотр своих записей}
\ucfield{Актор}{Клиент.}
\ucfield{Предусловия}{Клиент авторизован.}
\ucfield{Триггер}{Открыт раздел «Мои записи».}
\ucfield{Основной сценарий}{Система выбирает записи текущего клиента и показывает услугу, мастера, дату, начало и окончание, зафиксированную стоимость и статус. Будущие записи отделены от истории. Для доступных к отмене записей показано соответствующее действие.}
\ucfield{Альтернативы}{При отсутствии записей выводится пустой список с пояснением. Попытка получить чужую запись отклоняется.}
\ucfield{Постусловия}{Клиент получил сведения только о собственных записях; данные не изменены.}

\subsection{П5. Отмена своей записи}
\ucfield{Актор}{Клиент.}
\ucfield{Предусловия}{Клиент авторизован и выбрал собственную будущую действующую запись.}
\ucfield{Триггер}{Клиент подтверждает отмену.}
\ucfield{Основной сценарий}{Сервер проверяет владельца, блокирует соответствующего мастера и повторно проверяет статус и время начала. Затем переводит запись в статус «Отменена», сохраняет время и инициатора отмены. После фиксации интервал становится доступен при следующем расчёте.}
\ucfield{Альтернативы}{Отмена чужой или уже начавшейся записи запрещена. Для уже отменённой записи возвращается её актуальный статус без повторного изменения.}
\ucfield{Постусловия}{Запись сохранена в истории как отменённая и не блокирует время; при отказе состояние не изменилось.}

\subsection{П6. Управление услугами, мастерами и их связями}
\ucfield{Актор}{Администратор.}
\ucfield{Предусловия}{Выполнен вход с ролью администратора.}
\ucfield{Триггер}{Администратор создаёт или редактирует объект справочника.}
\ucfield{Основной сценарий}{Администратор вводит сведения об услуге или мастере и указывает допустимые связи. Система проверяет обязательные поля, положительную длительность, неотрицательную стоимость и уникальность пары «мастер--услуга». После проверки влияния на будущие записи сохраняются изменения.}
\ucfield{Альтернативы}{Некорректные данные отклоняются. Деактивация и удаление связи при соответствующих будущих действующих записях отклоняются с объяснением. Физическое удаление используемого объекта не предоставляется.}
\ucfield{Постусловия}{Справочники обновлены; условия ранее созданных записей сохранены.}

\subsection{П7. Управление рабочим временем}
\ucfield{Актор}{Администратор.}
\ucfield{Предусловия}{Выполнен административный вход; мастер существует.}
\ucfield{Триггер}{Администратор выбирает мастера и дату для изменения графика.}
\ucfield{Основной сценарий}{Администратор задаёт один или несколько рабочих интервалов. Сервер блокирует мастера, проверяет порядок границ, отсутствие перехода через полночь, пересечений интервалов и нарушений будущих записей. Система сохраняет допустимый график.}
\ucfield{Альтернативы}{Если изменение оставляет будущую запись вне графика, оно отклоняется с указанием конфликтующей записи. Удаление всех интервалов задаёт выходной и проходит ту же проверку.}
\ucfield{Постусловия}{Новый график участвует в следующем расчёте доступности; уже принятые записи остаются допустимыми.}

\subsection{П8. Работа с журналом записей}
\ucfield{Актор}{Администратор.}
\ucfield{Предусловия}{Выполнен административный вход.}
\ucfield{Триггер}{Администратор открывает журнал.}
\ucfield{Основной сценарий}{Система показывает записи с фильтрами по дате, мастеру, клиенту и статусу. Администратор открывает детали нужной записи.}
\ucfield{Альтернативный сценарий отмены}{Администратор выбирает будущую действующую запись, указывает причину и подтверждает отмену. Сервер в транзакции проверяет условия и сохраняет статус, причину и инициатора. Начавшуюся запись отменить нельзя; повторная отмена не меняет результат.}
\ucfield{Постусловия}{Получены сведения журнала; при отмене время освобождено, история сохранена.}

\subsection{П9. Проверка доступности перед сохранением}
\ucfield{Контекст}{Включаемый прецедент П3, самостоятельного внешнего актора нет.}
\ucfield{Предусловия}{Начата транзакция и получены необходимые блокировки.}
\ucfield{Основной сценарий}{Сервер проверяет активность услуги и мастера, связь между ними, допустимость даты и шага начала, актуальность условий, будущее время начала и полное вхождение в рабочий интервал. Затем проверяет отсутствие пересечений с действующими записями мастера.}
\ucfield{Альтернативы}{Нарушение любого условия возвращает конкретную причину отказа; сохранение не выполняется.}
\ucfield{Постусловия}{Выдано разрешение на сохранение в текущей транзакции либо отказ. Результат нельзя использовать после освобождения блокировки как гарантию доступности.}

\subsection{Связь требований и прецедентов}
\begin{center}
\begin{tabular}{ll}
\toprule Требования & Прецеденты \\ \midrule
ФТ-01 & П1 \\
ФТ-02, ФТ-03 & П2 \\
ФТ-04, ФТ-05 & П3, П9 \\
ФТ-06 & П4, П5 \\
ФТ-07 & П6 \\
ФТ-08 & П7 \\
ФТ-09 & П8 \\
ФТ-10 & П3, П5, П6, П8 \\
ФТ-11 & П1--П9 \\
\bottomrule
\end{tabular}
\end{center}

\section{Предлагаемая архитектура}
\subsection{Общий подход и технологии}
Предлагается трёхуровневая клиент-серверная архитектура с монолитной серверной частью. Она соответствует масштабу одного салона и позволяет выполнять операции записи в одной транзакции без распределённого взаимодействия.
\begin{itemize}
\item \textbf{Представление: Vue 3.} Одностраничный интерфейс с каталогом, выбором времени, личным разделом и административными формами. Сборка выполняется вне сервера демонстрации; в браузер передаются статические файлы.
\item \textbf{Прикладной уровень: Spring Boot и Spring MVC.} Spring Boot используется для настройки и запуска Java-приложения, Spring MVC --- для обработки HTTP-запросов. Это совместно используемые технологии, а не два альтернативных сервера.
\item \textbf{Данные: PostgreSQL.} Хранит пользователей, справочники, рабочие интервалы и записи; обеспечивает транзакции и блокировки.
\item \textbf{Вспомогательные средства:} Spring Security для аутентификации и прав доступа, Spring Data JPA для доступа к данным, Maven для сборки серверной части.
\end{itemize}
Точные версии Spring Boot, Java, средств сборки и библиотек фиксируются после проверки среды \texttt{helios} и согласования с преподавателем. Требования конкретной ветки Spring Boot сверяются с официальной документацией \cite{spring-req}. Возможность упаковки приложения в исполняемый JAR описана в официальном руководстве Spring \cite{spring-guide}.

\begin{figure}[htbp]
\centering
\begin{tikzpicture}[font=\small,
 box/.style={draw,rounded corners,align=center,text width=10cm,minimum height=1.1cm},
 link/.style={-{Stealth},thick}]
\node[box] (v) {Браузер клиента или администратора\\Интерфейс Vue};
\node[box,below=1cm of v] (c) {Приложение на \texttt{helios}\\Spring MVC: HTTP API и статические файлы Vue};
\node[box,below=0.8cm of c] (s) {Сервисы: расчёт времени, запись, график, справочники\\Spring Security: сессии и права доступа};
\node[box,below=0.8cm of s] (r) {Доступ к данным: Spring Data JPA\\Транзакции и согласованный протокол блокировок};
\node[box,below=0.8cm of r] (db) {PostgreSQL\\Постоянные данные системы};
\draw[link] (v)--node[right]{HTTP API: JSON; защищённый канал} (c);
\draw[link] (c)--(s);
\draw[link] (s)--(r);
\draw[link] (r)--node[right]{JDBC} (db);
\end{tikzpicture}
\caption{Логическая архитектура и основные связи}
\end{figure}

\subsection{Ответственность компонентов}
Vue отвечает за ввод, отображение и предварительную проверку формы. Расчёт доступности и решение о создании записи выполняются сервером. Контроллеры Spring MVC принимают запросы и возвращают ответы; сервисы реализуют бизнес-правила и определяют транзакционные границы; слой доступа к данным выполняет чтение, блокировки и сохранение.

Для базовой версии применяется серверная сессия. Административные операции доступны только соответствующей роли; при работе с личными записями проверяется владелец. Данные для интерфейса формируются отдельными объектами передачи данных, без хешей паролей и внутренних реквизитов БД.

Основные группы HTTP API: учётная запись и сессия; каталог услуг и мастеров; доступность; собственные записи; административные справочники, графики и журнал. Конкретные адреса и формат полей будут определены при проектировании интерфейсов. Конфликт времени отражается отдельной прикладной ошибкой с HTTP-статусом 409; интерфейс обновляет варианты и предлагает повторный выбор.

\section{Заключение}
В рамках первого этапа определены предметная область и границы информационной системы онлайн-записи в салон красоты, сформулированы её назначение, функциональные и нефункциональные требования, описаны роли, сущности и основные прецеденты. Определены требования к подбору свободного времени и предотвращению конфликтов записей.

Предлагаемая архитектура использует Vue, Spring Boot со Spring MVC и PostgreSQL и предусматривает демонстрацию серверной части на \texttt{helios}. Следующий организационный шаг --- согласование темы, прецедентов и технологий с преподавателем и проверка параметров серверной среды. После согласования можно переходить к детализации модели данных и реализации на следующих этапах курсовой работы.

\begin{thebibliography}{9}
\addcontentsline{toc}{section}{Список источников}
\bibitem{spring-req} Spring Boot. Системные требования. Официальная документация. URL: \url{https://docs.spring.io/spring-boot/system-requirements.html} (дата обращения: 15.09.2026).
\bibitem{spring-guide} Spring. Создание приложения Spring Boot. Официальное руководство. URL: \url{https://spring.io/guides/gs/spring-boot/} (дата обращения: 15.09.2026).
\end{thebibliography}
